\section{APORTES}

\subsection{PRÁCTICO}

El principal aporte práctico de la presente investigación consiste en la entrega de un sistema funcional capaz de procesar votos con una latencia constante de aproximadamente 165 milisegundos y un uso de CPU inferior al 7%, demostrando de manera empírica que la seguridad cuántica no resulta prohibitiva en términos de rendimiento cuando se implementa sobre hardware convencional. Este resultado se enmarca dentro de los rangos de viabilidad reportados por estudios previos sobre la evaluación de algoritmos post-cuánticos en dispositivos de recursos limitados (Lopez, Cadena y Rahman, 2025) y aporta evidencia adicional, específica al dominio de la votación electrónica, que no había sido reportada previamente con este nivel de detalle en la literatura consultada.\

Adicionalmente, el sistema desarrollado -estructurado en módulos independientes de criptografía, circuitos de conocimiento cero, contratos inteligentes, interfaz de línea de comandos y evaluación de rendimiento- constituye una base reutilizable para futuras investigaciones e implementaciones institucionales que requieran adoptar mecanismos de votación resistentes a la computación cuántica, sin necesidad de partir desde cero en el diseño de la arquitectura criptográfica.\

\subsection{TEÓRICO}

Desde el punto de vista teórico, el estudio propone un modelo de integración de funciones hash amigables para pruebas de conocimiento cero (Poseidon) con mecanismos de encapsulación de claves (KEM) post-cuánticos, orientado a generar compromisos (commitments) de voto que resulten seguros tanto frente a adversarios clásicos como frente a adversarios cuánticos. Este modelo extiende las propuestas previas de compromisos basados en retículos (Silde, 2021) al incorporar de manera explícita el secreto compartido derivado de ML-KEM como insumo del proceso de generación del commitment, lo cual no ha sido documentado de forma sistemática en la literatura revisada.\

El aporte teórico incluye, asimismo, la sistematización de un pipeline de votación E2E que articula de manera coherente los tres estándares FIPS 203, FIPS 204 y FIPS 205 (NIST, 2024) con el esquema de prueba Groth16 (Grassi et al., 2021; Tortola et al., 2024) y con los principios de verificabilidad E2E discutidos por Finogina y Herranz (2023), ofreciendo un marco de referencia que puede orientar el diseño de futuros sistemas de votación post-cuántica sobre infraestructuras de contratos inteligentes.\

